\editionPage[\pdfbookmark[0]{Corrigés des bilans}{bm-bilan-A-cor-auto-parent}]{bilan-A-cor-auto}{Bilan A · Corrigé}{Automatismes expliqués}{\hyperref[bilan-A-auto]{Énoncé p.\,\pageref*{bilan-A-auto}}}
\begin{blocEdition}{bilan-A-cor-auto}
\begin{questions}
\item $\res{2\,040\,000\,000}$.
\item $\res{3{,}507}$.
\item $\res{0{,}65<0{,}699<0{,}7}$.
\item $\res{14{,}39}$.
\item $25\times4=100$, deux décimales : $\res1$.
\item $\res{0{,}56}$.
\item $200=15\times13+5$ : quotient $\res{13}$, reste $\res5$.
\item $\square=\res{\frac57}$.
\item $\res{\frac12}$ et $\res{\frac34}$.
\item $18:3\times2=\res{12}$.
\item Le quart de 36 : $\res9$.
\item Deux fois plus : $\res9$ €.\end{questions}
\end{blocEdition}
\editionPage{bilan-A-cor-1}{Bilan A · Corrigé 1}{La médiathèque}{\hyperref[bilan-A-ex-1]{Énoncé p.\,\pageref*{bilan-A-ex-1}}}
\begin{blocEdition}{bilan-A-cor-1}
\editionEtape{Dans 4 ans}$12\,480+4\times350=\res{13\,880}$ livres.\editionEtape{Étagères}$12\,480=48\times260$ : $\res{260}$ étagères pleines, aucun livre restant.\editionEtape{Budget}$350\times18{,}50=\res{6\,475}$ €. Ordre de grandeur : $350\times20=7\,000$.
\end{blocEdition}
\editionPage{bilan-A-cor-2}{Bilan A · Corrigé 2}{Le marché}{\hyperref[bilan-A-ex-2]{Énoncé p.\,\pageref*{bilan-A-ex-2}}}
\begin{blocEdition}{bilan-A-cor-2}
\editionEtape{Clémentines}$2{,}5\times2{,}80=\res7$ €.\editionEtape{Cerises}250 g coûtent $6:3=2$ € ; 1 kg coûte $4\times2=\res8$ € le kg.\editionEtape{Monnaie}500 g de cerises coûtent 4 €. Total : $7+4=11$ € ; rendu : $\res9$ €.
\end{blocEdition}
\editionPage{bilan-A-cor-3}{Bilan A · Corrigé 3}{Le gâteau et les parts}{\hyperref[bilan-A-ex-3]{Énoncé p.\,\pageref*{bilan-A-ex-3}}}
\begin{blocEdition}{bilan-A-cor-3}
\editionEtape{Ensemble}$\frac3{12}+\frac4{12}=\res{\frac7{12}}$.\editionEtape{Reste}$1-\frac7{12}=\res{\frac5{12}}$.\editionEtape{Masse et pourcentage}$840:12\times5=\res{350}$ g. Léa a mangé $\frac14$, soit $\res{25\text{~\%}}$.
\end{blocEdition}
\editionPage{bilan-A-cor-4}{Bilan A · Corrigé 4}{Les cartes de collection}{\hyperref[bilan-A-ex-4]{Énoncé p.\,\pageref*{bilan-A-ex-4}}}
\begin{blocEdition}{bilan-A-cor-4}
\editionEtape{Barres}Quatre barres valent 96 : Sarah a $\res{24}$ cartes, Tom $\res{72}$.\editionEtape{Après le don}Tom : $72-12=\res{60}$ ; Sarah : $24+12=\res{36}$.\editionEtape{Égaliser}Au départ, l’écart vaut $72-24=48$ : Tom doit en donner la moitié, $\res{24}$. Chacun aurait alors 48 cartes.
\end{blocEdition}
\editionPage{bilan-B-cor-auto}{Bilan B · Corrigé}{Automatismes expliqués}{\hyperref[bilan-B-auto]{Énoncé p.\,\pageref*{bilan-B-auto}}}
\begin{blocEdition}{bilan-B-cor-auto}
\begin{questions}
\item $\res{3\,500}$ m.
\item $\res{4{,}5\ \mathrm{dm}^2}$.
\item $\res{20\ \mathrm{m}^2}$.
\item Environ $\res{62{,}8\ \mathrm{cm}}$.
\item $\res{2\text{ h }15\text{ min}}$.
\item $\res{70^\circ}$.
\item $\res{60^\circ}$.
\item $\res{145^\circ}$.
\item $\res{MB=6\ \mathrm{cm}}$.
\item $\res4$.
\item $\res{12\ \mathrm{cm}^3}$.
\item 3 et 6 : $\frac26=\res{\frac13}$.\end{questions}
\end{blocEdition}
\editionPage{bilan-B-cor-1}{Bilan B · Corrigé 1}{La cour du collège}{\hyperref[bilan-B-ex-1]{Énoncé p.\,\pageref*{bilan-B-ex-1}}}
\begin{blocEdition}{bilan-B-cor-1}
\editionEtape{Cour}Périmètre : $2\times(45+30)=\res{150\ \mathrm m}$. Aire : $45\times30=\res{1\,350\ \mathrm m^2}$.\editionEtape{Piste}$3{,}14\times20=\res{62{,}8\ \mathrm m}$ environ.\editionEtape{Course}$6\times62{,}8=\res{376{,}8\ \mathrm m}$. 3 min 30 s $=210$ s ; $210:6=\res{35\ \mathrm s}$ par tour.
\end{blocEdition}
\editionPage{bilan-B-cor-2}{Bilan B · Corrigé 2}{Le triangle et son cercle}{\hyperref[bilan-B-ex-2]{Énoncé p.\,\pageref*{bilan-B-ex-2}}}
\begin{blocEdition}{bilan-B-cor-2}
\editionEtape{Triangle}$\widehat{C}=180-40-60=\res{80^\circ}$.\editionEtape{Cercle}On trace deux médiatrices, par exemple celles de [AB] et [AC] ; leur point d’intersection O est à égale distance des trois sommets. Le cercle de centre O de rayon OA passe par A, B et C.\editionEtape{Symétrique}B et C sont sur l’axe : $B'=B$ et $C'=C$. La symétrie conserve les longueurs et les angles : $\res{A'B'=7\ \mathrm{cm}}$ et l’angle de sommet $B'$ mesure $\res{60^\circ}$.\label{fig-bilan-B-cor-2}\figureGeo[corrige]{bilan-triangle}

\end{blocEdition}
\editionPage{bilan-B-cor-3}{Bilan B · Corrigé 3}{La sortie au zoo}{\hyperref[bilan-B-ex-3]{Énoncé p.\,\pageref*{bilan-B-ex-3}}}
\begin{blocEdition}{bilan-B-cor-3}
\editionEtape{Trajet}$\res{1\text{ h }35\text{ min}}$.\editionEtape{Visite}$3{,}5$ h $=3$ h 30 min : fin à $\res{13\text{ h }45}$.\editionEtape{Allée}$9\times50=\res{450\ \mathrm m}$ ; $450:60=7{,}5$ min, soit $\res{7\text{ min }30\text{ s}}$.
\end{blocEdition}
\editionPage{bilan-B-cor-4}{Bilan B · Corrigé 4}{Les jeux de la kermesse}{\hyperref[bilan-B-ex-4]{Énoncé p.\,\pageref*{bilan-B-ex-4}}}
\begin{blocEdition}{bilan-B-cor-4}
\editionEtape{Urne}20 boules : $\frac6{20}=\frac{30}{100}=\res{30\text{~\%}}$.\editionEtape{Lots}La $\res{\text{loterie}}$, avec 42 lots ; total : $35+28+42+15=\res{120}$ lots.\editionEtape{Sans les jaunes}10 boules restent, dont 6 vertes : $\frac6{10}=\res{0{,}6}$, soit 60~\%.
\end{blocEdition}
